%% II. Tổng quan tài liệu và khoảng trống nghiên cứu. Được access_vi.tex \input.
\section{Tổng quan tài liệu}
\label{sec:litreview}

\subsection{Từ khác biệt nhóm đến động lực của một phiên}
Các nghiên cứu hình ảnh thần kinh đầu tiên về thiền so sánh giữa các nhóm: người thiền lâu năm có hoạt động DMN giảm khi thực hành và kết nối DMN thay đổi \cite{Brewer2011,Garrison2015}, và một phân tích gộp 78 nghiên cứu hình ảnh chức năng vẽ bản đồ các vùng được huy động bởi các kiểu thiền khác nhau \cite{Fox2016}. Các công trình sau chuyển sang động lực bên trong một phiên. Hasenkamp và cộng sự \cite{Hasenkamp2012} chia thực hành chú ý tập trung thành một chu kỳ gồm lang thang tâm trí, nhận ra lang thang, chuyển chú ý và duy trì tập trung; Lutz và cộng sự \cite{Lutz2008} mô tả chú ý tập trung và quan sát mở như hai chiến lược điều tiết chú ý; và Christoff và cộng sự \cite{Christoff2016} coi lang thang tâm trí là suy nghĩ tự phát được định hình bởi các ràng buộc có chủ đích và tự động. Các cách giải thích theo xử lý dự đoán mô tả thiền như một sự làm phẳng dần hệ thứ bậc dự đoán \cite{Laukkonen2021}, và một mô hình suy luận chủ động mô phỏng siêu nhận thức và kiểm soát chú ý trong chú ý tập trung \cite{SandvedSmith2021}. Các cách giải thích này mang tính định tính hoặc là mô hình sinh, và không cách nào nêu các tham số của nó ước lượng được chính xác đến đâu từ dữ liệu mà nghiên cứu thiền thu thập.

\subsection{Khung lý thuyết Bộ não như một giao diện truy vấn}
Trong các cách giải thích động lực, khung BQI của Vũ \cite{BQI} là khung viết động lực của một phiên thành một mô hình tường minh có tham số: nó nối hoạt động DMN, thông tin tích hợp và QoC do người tập báo cáo qua một hệ ODE có tốc độ và hệ số khuếch đại fit từ dữ liệu phân tích gộp, và đề xuất điều khiển chúng bằng một giao diện não--máy vòng kín \cite[Mục~7, 9 và 10]{BQI}. Phần tổng hợp tài liệu của BQI cho biết tín hiệu BOLD của DMN giảm gộp 39,4\,\% qua 16 nghiên cứu ($d=1{,}70$) và hoạt động DMN tương quan nghịch mạnh với QoC ($r=-0{,}89$) \cite[Mục~7.3--7.4]{BQI}. Trong bài đồng hành \cite{companion}, chúng tôi đã chỉ ra rằng phương trình QoC của BQI không có mức dừng và thay hệ bằng mô hình không gian trạng thái M1. BQI là nền tảng lý thuyết của nghiên cứu này; Mục~\ref{sec:theory} trình bày các phương trình của nó và dẫn ra M1 từ đó.

\subsection{Chỉ dấu EEG của hoạt động DMN và của tích hợp}
DMN được xác định từ chụp cắt lớp phát xạ positron và fMRI \cite{Raichle2001}, và các tương ứng EEG của nó là gián tiếp. Các nghiên cứu ghi đồng thời EEG--fMRI liên hệ dao động tự phát của công suất các dải tần với các mạng trạng thái nghỉ: công suất alpha đồng biến âm với một mạng chú ý trán--đỉnh và công suất beta 17--23~Hz đồng biến dương với các vùng của mạng mặc định \cite{Laufs2003}, và mỗi mạng trạng thái nghỉ có một tổ hợp nhịp EEG riêng, trong đó DMN gắn với công suất alpha và beta \cite{Mantini2007}. Vì vậy công suất alpha vùng sau cùng lắm là một đại diện gián tiếp của hoạt động DMN, và được dùng đúng như thế ở đây. Trong nghiên cứu thiền, một tổng quan hệ thống gắn chánh niệm với công suất alpha và theta tăng \cite{Lomas2015}; lang thang tâm trí khi tập trung vào hơi thở đi kèm công suất delta và theta nhiều hơn, alpha và beta ít hơn \cite{Braboszcz2011}; người tập lâu năm báo cáo ít lang thang tâm trí hơn, với các tương quan EEG của trạng thái được báo cáo \cite{Brandmeyer2018}; và khác biệt phổ giữa thiền và lang thang tâm trí không giống nhau ở người thiền lâu năm và người mới \cite{RodriguezLarios2021}. BQI coi công suất theta tăng là dấu hiệu phổ đặc trưng của thiền lâu năm \cite[Phụ lục~B.2]{BQI}. Về tích hợp, bản thân $\Phi$ không tính được cho hệ cỡ bộ não; chỉ số độ phức tạp nhiễu loạn nén đáp ứng của vỏ não với kích thích từ xuyên sọ bằng thuật toán Lempel--Ziv \cite{Casali2013}, độ phức tạp Lempel--Ziv của EEG tự phát giảm khi gây mê bằng propofol \cite{Schartner2015}, và số mũ aperiodic của phổ công suất là một phương án dựa trên phổ \cite{Donoghue2020}. BQI dùng chỉ số độ phức tạp làm đại diện cho $\Phi$ và báo cáo giá trị cao hơn ở người thiền lâu năm \cite[Mục~6.3]{BQI}.

\subsection{Khả năng xác định và thiết kế thí nghiệm cho mô hình động lực}
Liệu các tham số của một mô hình động lực có xác định được từ dữ liệu không nhiễu hay không là câu hỏi về khả năng xác định cấu trúc \cite{Bellman1970}; trong một mô hình tham số chính quy, một điểm xác định được cục bộ khi và chỉ khi ma trận thông tin không suy biến tại đó \cite{Rothenberg1971}. Khả năng xác định thực hành hỏi liệu dữ liệu hữu hạn, có nhiễu có xác định được tham số không, và hợp lý profile là công cụ chuẩn cho câu hỏi này \cite{Raue2009,Wieland2021}. Cận Cramér--Rao \cite{Rao1945,Cramer1946} biến thông tin Fisher thành cận dưới của phương sai của mọi ước lượng không chệch và là nền của thiết kế thí nghiệm tối ưu trong sinh học hệ thống \cite{Kreutz2009}. Với một quá trình ẩn tuyến tính--Gauss được lấy mẫu không đều, bộ lọc Kalman cho hàm hợp lý chính xác \cite{Kalman1960,Harvey1989,Durbin2012}. Khi mô hình sai, thông tin quan sát và thông tin kỳ vọng tách nhau, đó là cơ sở của phép chẩn đoán sai đặc tả dùng ở đây \cite{White1982}. Các công cụ này chưa được áp dụng cho thiết kế lấy mẫu trải nghiệm trong nghiên cứu thiền trước bài đồng hành \cite{companion}, vốn làm điều đó trên mô phỏng.

\subsection{Lấy mẫu trải nghiệm, đăng ký trước và lặp lại}
Lấy mẫu trải nghiệm bằng lần hỏi là cách chuẩn để đo lang thang tâm trí, và cách đặt câu hỏi và chọn thời điểm hỏi làm thay đổi tỉ lệ được báo cáo \cite{Weinstein2018}. Một dự án lặp lại quy mô lớn cho thấy chỉ khoảng một phần ba hiệu ứng tâm lý học đã công bố lặp lại được với kết quả có ý nghĩa \cite{OSC2015}, và đăng ký trước tách phân tích khẳng định khỏi phân tích khám phá bằng cách cố định kế hoạch trước khi xem dữ liệu \cite{Nosek2018}. Các bộ dữ liệu EEG lấy mẫu trải nghiệm mở \cite{Delorme2024ds,RodriguezLarios2021} nay cho phép kiểm định độc lập, có đăng ký trước, một mô hình mà không cần thu dữ liệu mới.

\subsection{Khoảng trống nghiên cứu}
\label{sec:gap}
Từ tổng quan trên có ba khoảng trống. Thứ nhất, chưa phân tích thiết kế nào cho một mô hình động lực của thiền được kiểm trên EEG lấy mẫu trải nghiệm thật; phân tích của bài đồng hành là \emph{in silico}. Thứ hai, các giá trị tham số của BQI được fit từ các tóm tắt phân tích gộp của tương phản fMRI và tương quan cắt ngang \cite[Mục~7.2--7.5]{BQI} và chưa bao giờ được ước lượng từ chuỗi thời gian EEG của từng phiên, trong khi chính các đại diện EEG của DMN và của $\Phi$ còn chưa được hiệu chuẩn. Thứ ba, chưa hằng số tốc độ nào của một mô hình thiền được lặp lại giữa các phòng thí nghiệm theo một kế hoạch đã đăng ký trước. Nghiên cứu này giải quyết ba khoảng trống đó bằng các câu hỏi nghiên cứu RQ1--RQ4.
